\section{Conditional readout uncertainty inside sample gaps}
The following bound separates uncertainty in the last Gaussian layer
from uncertainty in the hidden features. It applies to the actual
posterior and helps locate what remains missing from a sharper
population argument.

\begin{proposition}[Gaussian readout variance between observations]
Condition on the hidden weights $W=((a_j,b_j))_{j=1}^N$ and the sample.
Let $X_-<X_+$ be two observed inputs and let $x\in[X_-,X_+]$.
Write $h=X_+-X_-$ and
$Q_W=N^{-1}\sum_j(a_j^2+b_j^2)$.
Then
\[
 \operatorname{Var}_{\Pi}(z(x)\mid W,S)
 \le \frac{2n\lambda\nu}{N}+
       \frac{2\nu Q_Wh^2}{N}. \tag{RV1}
\]
For bounded labels, averaging over the hidden-weight posterior gives
\[
 \mathbb E_{W\mid S}\operatorname{Var}_{\Pi}(z(x)\mid W,S)
 \le\frac{2n\lambda\nu}{N}
 +\frac{2\nu h^2}{N}\left(3\nu+\frac1{2\lambda}\right). \tag{RV2}
\]
The same statement holds with $1$ replaced by $\widehat{Y^2}$ in
the last numerator when the observed labels are not bounded.
\end{proposition}
\begin{proof}
Let $H$ be the training feature matrix and let
$\sigma^2=n\lambda\nu/N$ be the actual observation variance.
The conditional Gaussian readout posterior has covariance
\[
 \Sigma=(N^2\nu^{-1}I+\sigma^{-2}H^\top H)^{-1}.
\]
Consequently $\Sigma\preceq(\nu/N^2)I$ and
$H\Sigma H^\top\preceq\sigma^2 I_n$.
Write $t=(x-X_-)/h$ and
\[
 h_W(x)=(1-t)h_W(X_-)+t h_W(X_+)+e(x).
\]
The ReLU Lipschitz property gives
$|e_j(x)|\le |a_j|h$.
The conditional variance of the first two terms is at most
$\sigma^2((1-t)^2+t^2)\le\sigma^2$, while that of the residual
is at most $(\nu/N^2)h^2\sum_j a_j^2\le\nu Q_Wh^2/N$.
Using $\operatorname{Var}(U+V)\le2\operatorname{Var}U+
2\operatorname{Var}V$ proves (RV1).

For completeness, Gaussian radial integration by parts gives
\[
 \mathbb E_\Pi Q_N
 =3\nu-\frac2\lambda\mathbb E_\Pi P_n[z(z-Y)]
 \le3\nu+\frac{\widehat{Y^2}}{2\lambda}.
\]
Indeed the Gaussian density is proportional to
$e^{-\omega Q_N/2}$, the parameter dimension is $3N$, and scaling
all parameters scales $z$ quadratically. Since $Q_W\le Q_N$,
averaging (RV1) proves (RV2).
\end{proof}

For a design with maximum adjacent sample gap
$h_{\max}=O(\log(n/\delta)/n)$ and $N\ge n$, this conditional
variance is at most
$O(\lambda+\log^2(n/\delta)/(n^3\lambda))$ in the sample range.
This does not control the variation of the conditional posterior
mean as the hidden weights vary. It also makes no assertion about
extrapolation beyond the smallest and largest observed inputs.
Consequently (RV2) alone does not bound the full posterior's
population squared error.
